\frfilename{mt224.tex}
\versiondate{29.9.04}
\copyrightdate{1997}

\def\chaptername{Teorema Dasar Kalkulus}
\def\sectionname{Fungsi bervariasi terbatas}

\newsection{224}

Sekarang saya beralih ke masalah kedua dari dua masalah yang dibahas dalam
bab ini: mengidentifikasi fungsi-fungsi real yang merupakan integral tak
tentu. Saya menggunakan kesempatan ini untuk memberikan pengantar singkat
tentang teori fungsi bervariasi terbatas, yang juga menarik sebagai objek kajian tersendiri
dan akan penting dalam Bab 28. Saya memberikan pencirian dasar fungsi-fungsi
ini sebagai selisih fungsi-fungsi monoton (224D), bersama contoh yang
mewakili sifat-sifat elementernya.

\leader{224A}{Definisi} Misalkan $f$ suatu fungsi bernilai real dan $D$
suatu subhimpunan $\Bbb R$. Saya mendefinisikan $\Var_D(f)$, yakni
{\bf variasi (total) dari $f$ pada $D$}, sebagai berikut. Jika
$D\cap\dom f=\emptyset$, maka $\Var_D(f)=0$. Jika tidak, $\Var_D(f)$ ialah

\Centerline{$\sup\{\sum_{i=1}^n|f(a_i)-f(a_{i-1})|:
  a_0,a_1,\ldots,a_n\in D\cap\dom f,\,a_0\le a_1\le\ldots\le a_n\}$,}

\noindent dengan membolehkan $\Var_{D}(f)=\infty$.
Jika $\Var_{D}(f)$ berhingga, kita mengatakan bahwa $f$ {\bf bervariasi
terbatas} pada $D$. \cmmnt{Jika konteksnya tampak jelas,} saya dapat
menuliskan $\Var f$ untuk $\Var_{\dom f}(f)$, dan mengatakan bahwa $f$
sekadar `{\bf bervariasi terbatas}' jika nilai ini berhingga.

\leader{224B}{Catatan-catatan}\cmmnt{ {\bf (a)} Dalam bab ini, perhatian
kita hampir seluruhnya tertuju pada kasus ketika $D$ suatu interval tertutup terbatas
yang termuat dalam $\dom f$. Rumusan umum tersebut akan berguna untuk beberapa
pertanyaan teknis yang muncul dalam Bab 28; tetapi jika itu membuat Anda lebih
nyaman, untuk sementara tidak ada yang hilang jika Anda mengandaikan bahwa
$D$ suatu interval.

\spheader 224Bb Jelas bahwa}

\Centerline{$\Var_D(f)=\Var_{D\cap\dom f}(f)=\Var(f\restr D)$}

\noindent untuk semua $D$, $f$.

\leader{224C}{Proposisi} (a) Jika $f$, $g$ dua fungsi bernilai real dan
$D\subseteq\Bbb R$, maka

\Centerline{$\Var_{D}(f+g)\le\Var_{D}(f)+\Var_{D}(g)$.}

(b) Jika $f$ suatu fungsi bernilai real, $D\subseteq\Bbb R$, dan
$c\in\Bbb R$, maka $\Var_{D}(cf)=|c|\Var_{D}(f)$.

(c) Jika $f$ suatu fungsi bernilai real, $D\subseteq\Bbb R$, dan
$x\in \Bbb R$, maka

\Centerline{$\Var_{D}(f)
\ge\Var_{D\cap\ocint{-\infty,x}}(f)+\Var_{D\cap\coint{x,\infty}}(f)$,}

\noindent dengan kesamaan jika $x\in D\cap\dom f$.

(d) Jika $f$ suatu fungsi bernilai real dan
$D\subseteq D'\subseteq\Bbb R$, maka $\Var_D(f)\le\Var_{D'}(f)$.

(e) Jika $f$ suatu fungsi bernilai real dan $D\subseteq\Bbb R$, maka
$|f(x)-f(y)|\le\Var_D(f)$ untuk semua $x$, $y\in D\cap\dom f$; jadi jika
$f$ bervariasi terbatas pada $D$, maka $f$ terbatas pada $D\cap\dom f$ dan
(jika $D\cap\dom f\ne\emptyset$)

\Centerline{$\sup_{y\in D\cap\dom f}|f(y)|\le |f(x)|+\Var_D(f)$}

\noindent untuk setiap $x\in D\cap\dom f$.

(f) Jika $f$ suatu fungsi bernilai real monoton dan $D\subseteq\Bbb R$
beririsan dengan $\dom f$, maka
\discrcenter{468pt}
{$\Var_D(f)=\sup_{x\in D\cap\dom f}f(x)-\inf_{x\in D\cap\dom f}f(x)$.}

\proof{{\bf (a)} Jika $D\cap\dom(f+g)=\emptyset$, pernyataan ini sepele,
karena $\Var_D(f)$ dan $\Var_D(g)$ tentu nonnegatif. Jika tidak, apabila
$a_0\le\ldots\le a_n$ dalam $D\cap\dom(f+g)$, maka

$$\eqalign{\sum_{i=1}^n|(f+g)(a_i)-(f+g)(a_{i-1})|
&\le \sum_{i=1}^n|f(a_i)-f(a_{i-1})|
+\sum_{i=1}^n|g(a_i)-g(a_{i-1})|\cr
&\le\Var_{D}(f)+\Var_{D}(g);\cr}$$

\noindent karena $a_0,\ldots,a_n$ sebarang,
$\Var_{D}(f+g)\le\Var_{D}(f)+\Var_{D}(g)$.

\medskip

{\bf (b)}

\Centerline{$\sum_{i=1}^n|(cf)(a_i)-(cf)(a_{i-1})|
=|c|\sum_{i=1}^n|f(a_i)-f(a_{i-1})|$}

\noindent kapan pun
$a_0\le\ldots\le a_n$ dalam $D\cap\dom f$.

\medskip

{\bf (c)(i)} Jika salah satu dari
$D\cap\ocint{-\infty,x}\cap\dom f$ dan
$D\cap\coint{x,\infty}\cap\dom f$ kosong, pernyataan ini sepele. Jika
$a_0\le\ldots\le a_m$ dalam $D\cap\ocint{-\infty,x}\cap\dom f$ dan
$b_0\le\ldots\le b_n$ dalam $D\cap\coint{x,\infty}\cap\dom f$, maka

$$\eqalign{\sum_{i=1}^m|f(a_i)-f(a_{i-1})|
+\sum_{j=1}^n|f(b_j)-f(b_{j-1})|
&\le\sum_{i=1}^{m+n+1}|f(a_i)-f(a_{i-1})|\cr
&\le\Var_{D}(f),\cr}$$

\noindent jika kita menuliskan $a_i=b_{i-m-1}$ untuk
$m+1\le i\le m+n+1$. Jadi

\Centerline{$\Var_{D\cap\ocint{-\infty,x}}(f)
+\Var_{D\cap\coint{x,\infty}}(f)
\le\Var_D(f)$.}

\medskip

\quad{\bf (ii)} Sekarang andaikan bahwa $x\in D\cap\dom f$.
Jika $a_0\le\ldots\le a_n$ dalam $D\cap\dom f$, dan
$a_0\le x\le a_n$, misalkan $k$ sedemikian sehingga
$x\in [a_{k-1},a_k]$; maka

$$\eqalign{\sum_{i=1}^n|f(a_i)-f(a_{i-1})|
&\le\sum_{i=1}^{k-1}|f(a_i)-f(a_{i-1})|
+|f(x)-f(a_{k-1})|\cr
&\mskip150mu +|f(a_k)-f(x)|+\sum_{i=k+1}^n|f(a_i)-f(a_{i-1})|\cr
&\le\Var_{D\cap\ocint{-\infty,x}}(f)
+\Var_{D\cap\coint{x,\infty}}(f)
\cr}$$

\noindent (dengan menghitung jumlah kosong $\sum_{i=1}^0$ dan
$\sum_{i=n+1}^n$ sebagai $0$). Jika $x\le a_0$, maka
$\sum_{i=1}^n|f(a_i)-f(a_{i-1})|\le\Var_{D\cap\coint{x,\infty}}(f)$;
jika $x\ge a_n$, maka
$\sum_{i=1}^n|f(a_i)-f(a_{i-1})|\le\Var_{D\cap\ocint{-\infty,x}}(f)$.
Jadi

\Centerline{$\sum_{i=1}^n|f(a_i)-f(a_{i-1})|
\le\Var_{D\cap\ocint{-\infty,x}}(f)
+\Var_{D\cap\coint{x,\infty}}(f)
$}

\noindent dalam semua kasus; karena $a_0,\ldots,a_n$ sebarang,

\Centerline{$\Var_D(f)
\le\Var_{D\cap\ocint{-\infty,x}}(f)
+\Var_{D\cap\coint{x,\infty}}(f)$.}

\noindent Jadi kedua ruas sama.

\medskip

{\bf (d)} sepele.

\medskip

{\bf (e)} Jika $x$, $y\in D\cap\dom f$ dan $x\le y$, maka

\Centerline{$|f(x)-f(y)|=|f(y)-f(x)|\le \Var_D(f)$}

\noindent menurut definisi $\Var_D$; dan hal yang sama benar jika
$y\le x$. Jadi tentu saja $|f(y)|\le|f(x)|+\Var_D(f)$.

\medskip

{\bf (f)} Jika $f$ tak-menurun, maka

$$\eqalign{\Var_D(f)
&=\sup\{\sum_{i=1}^n|f(a_i)-f(a_{i-1})|:a_0,a_1,\ldots,a_n\in
D\cap\dom f,\,a_0\le
a_1\le\ldots\le a_n\}\cr
&=\sup\{\sum_{i=1}^nf(a_i)-f(a_{i-1}):a_0,a_1,\ldots,a_n\in
D\cap\dom f,\,a_0\le
a_1\le\ldots\le a_n\}\cr
&=\sup\{f(b)-f(a):a,\,b\in D\cap\dom f,\,a\le b\}\cr
&=\sup_{b\in D\cap\dom f}f(b)-\inf_{a\in D\cap\dom f}f(a).\cr}$$

\noindent Jika $f$ tak-menaik, maka

$$\eqalign{\Var_D(f)
&=\sup\{\sum_{i=1}^n|f(a_i)-f(a_{i-1})|:a_0,a_1,\ldots,a_n\in
D\cap\dom f,\,a_0\le
a_1\le\ldots\le a_n\}\cr
&=\sup\{\sum_{i=1}^nf(a_{i-1})-f(a_i):a_0,a_1,\ldots,a_n\in
D\cap\dom f,\,a_0\le
a_1\le\ldots\le a_n\}\cr
&=\sup\{f(a)-f(b):a,\,b\in D\cap\dom f,\,a\le b\}\cr
&=\sup_{a\in D\cap\dom f}f(a)-\inf_{b\in D\cap\dom f}f(b).\cr}$$
}%end of proof of 224C
\leader{224D}{Teorema} Untuk setiap fungsi bernilai real $f$ dan setiap
himpunan $D\subseteq \Bbb R$, pernyataan-pernyataan berikut ekuivalen:

\quad (i) terdapat dua fungsi terbatas dan tak-menurun $f_1$,
$f_2:\Bbb R\to\Bbb R$ sedemikian sehingga
$f=f_1-f_2$ pada $D\cap\dom f$;

\quad (ii) $f$ bervariasi terbatas pada $D$;

\quad (iii) terdapat fungsi terbatas dan tak-menurun $f_1$,
$f_2:\Bbb R\to\Bbb R$ sedemikian sehingga
$f=f_1-f_2$ pada $D\cap\dom f$ dan $\Var_D(f)=\Var f_1+\Var f_2$.

\proof{{\bf (i)$\Rightarrow$(ii)} Jika $f:\Bbb R\to\Bbb R$ terbatas
dan tak-menurun, maka
$\Var f=\sup_{x\in\Bbb R}f(x)-\inf_{x\in\Bbb R}f(x)$ berhingga. Jadi,
jika $f$ pada $D\cap\dom f$ berimpit dengan $f_1-f_2$, dengan $f_1$ dan
$f_2$ terbatas dan tak-menurun, maka

$$\eqalign{\Var_D(f)
&=\Var_{D\cap\dom f}(f)
\le\Var_{D\cap\dom f}(f_1)+\Var_{D\cap\dom f}(f_2)\cr
&\le\Var f_1+\Var f_2
<\infty,}$$

\noindent dengan menggunakan (a), (b), dan (d) dari 224C.

\medskip

{\bf(ii)$\Rightarrow$(iii)} Andaikan bahwa $f$ bervariasi terbatas pada
$D$. Tetapkan $D'=D\cap\dom f$. Jika $D'=\emptyset$, kita dapat mengambil
kedua $f_j$ sebagai fungsi nol; jadi, mulai sekarang andaikan bahwa
$D'\ne\emptyset$. Tuliskan

\Centerline{$g(x)=\Var_{D\cap\ocint{-\infty,x}}(f)$}

\noindent untuk $x\in D'$. Maka $g_1=g+f$ dan $g_2=g-f$ keduanya
tak-menurun. \Prf\ Jika $a$, $b\in D'$ dan $a\le b$, maka

\Centerline{$g(b)=g(a)+\Var_{D\cap[a,b]}(f)\ge g(a)+|f(b)-f(a)|$.}

\noindent Jadi

\Centerline{$g_1(b)-g_1(a)=g(b)-g(a)+f(b)-f(a)$,
\quad$g_2(b)-g_2(a)=g(b)-g(a)-f(b)+f(a)$}

\noindent keduanya nonnegatif.\ \Qed

Sekarang terdapat fungsi-fungsi tak-menurun
$h_1$, $h_2:\Bbb R\to\Bbb R$ yang masing-masing memperluas $g_1$, $g_2$,
sedemikian sehingga $\Var h_j=\Var g_j$ untuk masing-masing $j$. \Prf\ $f$
terbatas pada $D$, menurut 224Ce, dan $g$ terbatas semata-mata karena
$\Var_D(f)<\infty$, sehingga $g_j$ terbatas. Tetapkan
$c_j=\inf_{x\in D'}g_j(x)$ dan

\Centerline{$h_j(x)=\sup(\{c_j\}\cup\{g_j(y):y\in D',\,y\le x\})$}

\noindent untuk setiap $x\in\Bbb R$; definisi ini memenuhi syarat.\ \QeD\
Perhatikan bahwa untuk $x\in D'$,

\Centerline{$h_1(x)+h_2(x)=g_1(x)+g_2(x)=g(x)+f(x)+g(x)-f(x)=2g(x)$,}

\Centerline{$h_1(x)-h_2(x)=2f(x)$.}

Sekarang, karena $g_1$ dan $g_2$ tak-menurun,

\Centerline{$\sup_{x\in D'}g_1(x)+\sup_{x\in D'}g_2(x)=\sup_{x\in
D'}g_1(x)+g_2(x)=2\sup_{x\in D'}g(x)$,}

\Centerline{$\inf_{x\in D'}g_1(x)+\inf_{x\in D'}g_2(x)=\inf_{x\in
D'}g_1(x)+g_2(x)=2\inf_{x\in D'}g(x)\ge 0$.}

\noindent Namun, ini berarti bahwa

\Centerline{$\Var h_1+\Var h_2
=\Var g_1+\Var g_2
=2\Var g
\le 2\Var_D(f)$,}

\noindent dengan menggunakan 224Cf tiga kali. Jadi, jika kita menetapkan
$f_j(x)=\bover12h_j(x)$
untuk $j\in\{1,2\}$ dan $x\in\Bbb R$, kita akan memperoleh fungsi-fungsi
tak-menurun sedemikian sehingga

\Centerline{$f_1(x)-f_2(x)=f(x)$ untuk $x\in D'$,
\quad$\Var f_1+\Var f_2=\Bover12\Var h_1+\Bover12\Var h_2\le\Var_D(f)$.}

\noindent Karena tentu kita juga mempunyai

\Centerline{$\Var_D(f)\le\Var_D(f_1)+\Var_D(f_2)\le\Var f_1+\Var f_2$,}

\noindent kita melihat bahwa $\Var_D(f)=\Var f_1+\Var f_2$, dan (iii)
benar.

\medskip

{\bf (iii)$\Rightarrow$(i)} langsung.
}%end of proof of 224D

\leader{224E}{Korolari} Misalkan $f$ suatu fungsi bernilai real dan $D$
sembarang subhimpunan dari $\Bbb R$. Jika $f$ bervariasi terbatas pada $D$,
maka

\Centerline{$\lim_{x\downarrow a}\Var_{D\cap\ocint{a,x}}(f)
=\lim_{x\uparrow a}\Var_{D\cap\coint{x,a}}(f)=0$}
\noindent untuk setiap $a\in\Bbb R$, dan

\Centerline{$\lim_{a\to-\infty}\Var_{D\cap\ocint{-\infty,a}}(f)
=\lim_{a\to\infty}\Var_{D\cap\coint{a,\infty}}(f)=0$.}

\proof{{\bf (a)} Tinjau terlebih dahulu kasus ketika
$D=\dom f=\Bbb R$ dan $f$ merupakan fungsi terbatas dan tak-menurun. Maka

\Centerline{$\Var_{D\cap\ocint{a,x}}(f)
=\sup_{y\in\ocint{a,x}}f(x)-f(y)
=f(x)-\inf_{y>a}f(y)
=f(x)-\lim_{y\downarrow a}f(y)$,}

\noindent sehingga tentu saja

\Centerline{$\lim_{x\downarrow a}\Var_{D\cap\ocint{a,x}}(f)
=\lim_{x\downarrow a}f(x)-\lim_{y\downarrow a}f(y)=0$.}

\noindent Dengan cara yang sama,

\Centerline{$\lim_{x\uparrow a}\Var_{D\cap\coint{x,a}}(f)
=\lim_{y\uparrow a}f(y)-\lim_{x\uparrow a}f(x)=0$,}

\Centerline{$\lim_{a\to-\infty}\Var_{D\cap\ocint{-\infty,a}}(f)
=\lim_{a\to-\infty}f(a)-\lim_{y\to-\infty}f(y)=0$,}

\Centerline{$\lim_{a\to\infty}\Var_{D\cap\coint{a,\infty}}(f)
=\lim_{y\to\infty}f(y)-\lim_{a\to\infty}f(a)=0$.}

\medskip

{\bf (b)} Untuk kasus umum, definisikan $f_1$, $f_2$ dari $f$ dan $D$
seperti dalam 224D. Maka, untuk setiap interval $I$, kita mempunyai

\Centerline{$\Var_{D\cap I}(f)\le\Var_I(f_1)+\Var_I(f_2)$,}

\noindent sehingga hasil-hasil untuk $f$ mengikuti hasil-hasil untuk
$f_1$ dan $f_2$ yang telah dibuktikan dalam bagian (a).
}%end of proof of 224E

\leader{224F}{Korolari} Misalkan $f$ suatu fungsi bernilai real yang
bervariasi terbatas pada $[a,b]$, dengan $a<b$. Jika $\dom f$ beririsan
dengan setiap interval $\ocint{a,a+\delta}$ untuk $\delta>0$, maka

\Centerline{$\lim_{t\in\dom f,t\downarrow a}f(t)$}

\noindent ada dalam $\Bbb R$. Jika $\dom f$ beririsan dengan
$\coint{b-\delta,b}$ untuk setiap $\delta>0$, maka

\Centerline{$\lim_{t\in\dom f,t\uparrow b}f(t)$}

\noindent ada dalam $\Bbb R$.

\wheader{224F}{0}{0}{0}{36pt}

\proof{ Misalkan $f_1$, $f_2:\Bbb R\to\Bbb R$ fungsi-fungsi tak-menurun
sedemikian sehingga $f=f_1-f_2$ pada $[a,b]\cap\dom f$. Maka

\Centerline{$\lim_{t\in\dom f,t\downarrow a}f(t)
=\lim_{t\downarrow a}f_1(t)-\lim_{t\downarrow a}f_2(t)
=\inf_{t>a}f_1(t)-\inf_{t>a}f_2(t)$,}

\Centerline{$\lim_{t\in\dom f,t\uparrow b}f(t)
=\lim_{t\uparrow b}f_1(t)-\lim_{t\uparrow b}f_2(t)
=\sup_{t<b}f_1(t)-\sup_{t<b}f_2(t)$.}
}%end of proof of 224F

\leader{224G}{Korolari} Misalkan $f$, $g$ fungsi-fungsi bernilai real dan
$D$ suatu subhimpunan dari $\Bbb R$. Jika $f$ dan $g$ bervariasi terbatas
pada $D$, demikian pula $f\times g$.

\proof{{\bf (a)} Pokoknya adalah bahwa terdapat fungsi-fungsi terbatas,
tak-menurun, dan {\it nonnegatif} $f_1$, $f_2:\Bbb R\to\Bbb R$ sedemikian
sehingga $f=f_1-f_2$ pada $D\cap\dom f$. \Prf\ Kita tahu bahwa terdapat
fungsi terbatas dan tak-menurun $h_1$, $h_2$ sedemikian sehingga
$f=h_1-h_2$ pada $D\cap\dom f$. Tetapkan
$\gamma_i=\inf_{x\in\Bbb R}h_i(x)$ untuk $i=1$, $2$,

\Centerline{$\beta_1=\max(\gamma_1-\gamma_2,0)$,
\quad$\beta_2=\max(\gamma_2-\gamma_1,0)$,}

\Centerline{$f_1=h_1-\gamma_1+\beta_1$,
\quad$f_2=h_2-\gamma_2+\beta_2$;}

\noindent definisi ini memenuhi syarat.\ \Qed

\medskip

{\bf (b)} Sekarang, dengan mengambil fungsi-fungsi serupa $g_1$, $g_2$
sedemikian sehingga $g=g_1-g_2$ pada $D\cap\dom g$, kita mempunyai

\Centerline{$f\times g=f_1\times g_1-f_2\times g_1-f_1\times g_2
+f_2\times g_2$}

\noindent di setiap titik dalam
$D\cap\dom(f\times g)=D\cap\dom f\cap\dom g$; tetapi semua
$f_i\times g_j$ merupakan fungsi terbatas dan tak-menurun, sehingga
bervariasi terbatas, dan $f\times g$ pasti bervariasi terbatas pada $D$.
}%end of proof of 224G

\leader{224H}{Proposisi} Misalkan $f:D\to\Bbb R$ suatu fungsi bervariasi
terbatas, dengan $D\subseteq\Bbb R$. Maka $f$ kontinu pada $D$, kecuali
pada terhitung banyak titik.

\proof{ Untuk $n\ge 1$, tetapkan

$$\eqalign{A_n
&=\{x:x\in D,\text{ untuk setiap }\delta>0\text{ terdapat }y\in
D\cap[x-\delta,x+\delta]\cr
&\qquad\qquad\qquad\qquad\text{ sedemikian sehingga }|f(y)-f(x)|
  \ge\Bover1n\}.\cr}$$

\noindent Maka $\#(A_n)\le n\Var f$. \Prf\Quer\ Jika tidak, kita dapat
menemukan $x_0,\ldots,x_k\in A_n$ yang berbeda-beda dengan
$k+1>n\Var f$. Urutkan titik-titik ini sedemikian sehingga
$x_0<x_1<\ldots<x_k$. Tetapkan
$\delta=\bover12\min_{1\le i\le k}x_i-x_{i-1}>0$. Untuk setiap $i$,
terdapat $y_i\in D\cap[x_i-\delta,x_i+\delta]$ sedemikian sehingga
$|f(y_i)-f(x_i)|\ge\bover1n$. Ambil $x'_i$, $y'_i$ sebagai $x_i$, $y_i$
menurut urutan besarnya, sehingga $x'_i<y'_i$. Sekarang

\Centerline{$x'_0\le y'_0\le x'_1\le y'_1\le\ldots\le x'_k\le y'_k$,}
\noindent dan

\Centerline{$\Var f
\ge\sum_{i=0}^k|f(y'_i)-f(x'_i)|
=\sum_{i=0}^k|f(y_i)-f(x_i)|
\ge \Bover1n(k+1)
>\Var f$,}

\noindent yang mustahil.\ \Bang\Qed

Akibatnya, $A=\bigcup_{n\ge 1}A_n$ terhitung, karena merupakan
gabungan terhitung dari himpunan-himpunan berhingga. Namun, $A$ tepat
merupakan himpunan titik-titik dalam $D$ tempat $f$ tidak kontinu.
}%end of proof of 224H

\leader{224I}{Teorema} Misalkan $I\subseteq\Bbb R$ suatu interval, dan
$f:I\to\Bbb R$ suatu fungsi bervariasi terbatas. Maka $f$ dapat
didiferensialkan hampir di mana-mana dalam $I$, dan $f'$ terintegralkan
pada $I$, dengan

\Centerline{$\int_I|f'|\le\Var_I(f)$.}

\proof{{\bf (a)} Misalkan $f_1$ dan $f_2$ fungsi-fungsi tak-menurun
sedemikian sehingga $f=f_1-f_2$ di setiap titik dalam $I$ (224D). Maka
$f_1$ dan $f_2$ terdiferensialkan hampir di mana-mana (222A). Pada
setiap titik dalam $I$ selain mungkin titik-titik ujungnya, jika ada, $f$
terdiferensialkan jika $f_1$ dan $f_2$ terdiferensialkan, sehingga
$f'(x)$ terdefinisi untuk hampir setiap $x\in I$.

\medskip

{\bf (b)} Tetapkan $F(x)=\Var_{I\cap\ocint{-\infty,x}}f$ untuk
$x\in\Bbb R$. Jika $x$, $y\in I$ dan $x\le y$, maka

\Centerline{$F(y)-F(x)=\Var_{[x,y]}f\ge|f(y)-f(x)|$,}

\noindent menurut 224Cc; sehingga $F'(x)\ge|f'(x)|$ kapan pun $x$
merupakan titik interior dari $I$ dan kedua turunan itu ada, yang berlaku
hampir di mana-mana. Jadi $\int_I|f'|\le\int_IF'$. Namun, jika $a$,
$b\in I$ dan $a\le b$,

\Centerline{$\int_a^bF'\le F(b)-F(a)\le F(b)\le\Var f$.}

\noindent Sekarang $I$ dapat dinyatakan sebagai
$\bigcup_{n\in\Bbb N}[a_n,b_n]$, dengan
$a_{n+1}\le a_n\le b_n\le b_{n+1}$ untuk setiap $n$. Jadi

$$\eqalignno{\int_I|f'|
&\le\int_IF'
=\int F'\times\chi I\cr
&=\int\sup_{n\in\Bbb N}F'\times\chi[a_n,b_n]
=\sup_{n\in\Bbb N}\int F'\times\chi[a_n,b_n]\cr
\noalign{\noindent (menurut teorema B. Levi)}
&=\sup_{n\in\Bbb N}\int_{a_n}^{b_n}F'
\le\Var_I(f).\cr}$$
}%end of proof of 224I

\leader{224J}{}\cmmnt{ Hasil berikut tidak diperlukan dalam bab ini,
tetapi merupakan salah satu sifat yang paling berguna dari fungsi
bervariasi terbatas, dan akan digunakan berulang kali dalam Bab 28.

\medskip

\noindent}{\bf Proposisi} Misalkan $f$, $g$ fungsi-fungsi bernilai real
yang didefinisikan pada subhimpunan-subhimpunan $\Bbb R$, dan andaikan
bahwa $g$ terintegralkan pada interval $[a,b]$, dengan $a<b$, sedangkan
$f$ bervariasi terbatas pada $\ooint{a,b}$ dan didefinisikan hampir di
mana-mana dalam $\ooint{a,b}$. Maka $f\times g$ terintegralkan pada
$[a,b]$, dan

$$\bigl|\int_a^bf\times g\bigr|
\le\bigl(\lim_{x\in\dom f,x\uparrow b}|f(x)|+\Var_{\ooint{a,b}}(f)\bigr)
  \sup_{c\in[a,b]}\bigl|\int_a^cg\bigr|.$$

\proof{{\bf (a)} Menurut 224F,
$l=\lim_{x\in\dom f,x\uparrow b}f(x)$ terdefinisi. Tuliskan
$M=|l|+\Var_{\ooint{a,b}}(f)$. Perhatikan bahwa jika $y$ adalah sembarang
titik dari $\dom f\cap\ooint{a,b}$,

\Centerline{$|f(y)|\le|f(x)|+|f(x)-f(y)|
\le|f(x)|+\Var_{\ooint{a,b}}(f)\to M$}

\noindent ketika $x\uparrow b$ di dalam $\dom f$, sehingga
$|f(y)|\le M$. Selain itu, $f$ terukur pada $\ooint{a,b}$, karena
terdapat fungsi-fungsi monoton terbatas
$f_1$, $f_2:\Bbb R\to\Bbb R$ sedemikian sehingga $f=f_1-f_2$
di setiap titik dalam $\ooint{a,b}\cap\dom f$. Jadi $f\times g$ terukur,
didominasi oleh $M|g|$, dan terintegralkan pada $\ooint{a,b}$ maupun
$[a,b]$.

\medskip

{\bf (b)} Untuk $n\in\Bbb N$, $k\le 2^n$, tetapkan
$a_{nk}=a+2^{-n}k(b-a)$, dan untuk $1\le k\le 2^n$, pilih
$x_{nk}\in\dom f\cap\ooint{a_{n,k-1},a_{nk}}$. Definisikan
$f_n:\ocint{a,b}\to\Bbb R$ dengan menetapkan $f_n(x)=f(x_{nk})$ jika
$1\le k\le 2^n$ dan $x\in\ocint{a_{n,k-1},a_{nk}}$. Maka
$f(x)=\lim_{n\to\infty}f_n(x)$ kapan pun
$x\in\ooint{a,b}\cap\dom f$ dan $f$ kontinu di $x$, yang pasti berlaku
hampir di mana-mana (224H). Selanjutnya, perhatikan bahwa semua $f_n$
terukur, dan nilai mutlaknya dibatasi secara seragam oleh $M$. Jadi
$\{f_n\times g:n\in\Bbb N\}$ didominasi oleh fungsi terintegralkan
$M|g|$, dan Teorema Konvergensi Terdominasi Lebesgue menyatakan bahwa

\Centerline{$\int_a^bf\times g=\lim_{n\to\infty}\int_a^bf_n\times g$.}

\medskip

{\bf (c)} Untuk sementara, tetapkan $n\in\Bbb N$. Tetapkan
$K=\sup_{c\in[a,b]}|\int_a^cg|$. (Perhatikan bahwa $K$ berhingga karena
$c\mapsto\int_a^cg$ kontinu.) Maka

$$\eqalignno{\bigl|\int_a^bf_n\times g\bigr|
&=\bigl|\sum_{k=1}^{2^n}\int_{a_{n,k-1}}^{a_{nk}}f_n\times g\bigr|\cr
&=\bigl|\sum_{k=1}^{2^n}
    f(x_{nk})(\int_a^{a_{nk}}g-\int_a^{a_{n,k-1}}g)\bigr|\cr
&=\bigl|\sum_{k=1}^{2^n-1}
   (f(x_{nk})-f(x_{n,k+1}))\int_a^{a_{nk}}g
   +f(x_{n,2^n})\int_a^bg\bigr|\cr
&\le\bigl|f(x_{n,2^n})\bigr|\bigl|\int_a^bg\bigr|+\sum_{k=1}^{2^n-1}
   \bigl|f(x_{n,k+1})-f(x_{nk})\bigr|\bigl|\int_a^{a_{nk}}g\bigr|\cr
&\le(|f(x_{n,2^n})|+\Var_{\ooint{a,b}}(f))K
\to MK\cr}$$

\noindent ketika $n\to\infty$.

\medskip

{\bf (d)} Sekarang

\Centerline{$|\int_a^bf\times g|
=\lim_{n\to\infty}|\int_a^bf_n\times g|\le MK$,}

\noindent sebagaimana diperlukan.
}%end of proof of 224J

\leader{224K}{Fungsi bernilai kompleks}\cmmnt{ Sejauh ini saya telah
menganggap semua fungsi bernilai real. Ini memadai bagi keperluan bab ini,
tetapi dalam Bab 28 kita perlu meninjau fungsi bernilai kompleks yang
bervariasi terbatas, dan barangkali saya perlu menguraikan penyesuaian
(elementer) yang terlibat dalam perluasan ke kasus kompleks.
}

\spheader 224Ka Misalkan $D$ suatu subhimpunan dari $\Bbb R$ dan $f$
suatu fungsi bernilai kompleks. {\bf Variasi} $f$ pada $D$, $\Var_D(f)$,
adalah nol jika $D\cap\dom f=\emptyset$, dan jika tidak, variasi tersebut
adalah

\Centerline{$\sup\{\sum_{j=1}^n|f(a_j)-f(a_{j-1})|:
 a_0\le a_1\le\ldots\le a_n$ dalam $D\cap\dom f\}$,}

\noindent dengan mengizinkan nilai $\infty$.
Jika $\Var_{D}(f)$ berhingga, kita mengatakan bahwa $f$ {\bf bervariasi
terbatas} pada $D$.

\spheader 224Kb \dvro{Suatu}{Seperti dalam kasus real, suatu} fungsi
bernilai kompleks yang bervariasi terbatas pasti terbatas, dan
\Centerline{$\Var_{D}(f+g)\le\Var_{D}(f)+\Var_{D}(g)$,}

\Centerline{$\Var_{D}(cf)=|c|\Var_{D}(f)$,}

\Centerline{$\Var_{D}(f)\ge\Var_{D\cap\ocint{-\infty,x}}(f)
+\Var_{D\cap\coint{x,\infty}}(f)$}

\noindent untuk setiap
$x\in\Bbb R$, dengan kesamaan jika $x\in D\cap\dom f$,

\Centerline{$\Var_D(f)\le\Var_{D'}(f)$ kapan pun $D\subseteq
D'$\dvro{.}{;}}

\cmmnt{\noindent argumen-argumen 224C berlaku tanpa perubahan.}

\spheader 224Kc Suatu fungsi bernilai kompleks bervariasi terbatas jika
dan hanya jika bagian real dan bagian imajinernya sama-sama bervariasi
terbatas\dvro{. }{ (karena

\Centerline{$\max(\Var_{D}(\Real f),\Var_{D}(\Imag f))
\le\Var_{D}(f)\le\Var_{D}(\Real f)+\Var_{D}(\Imag f)$.)}

\noindent}Jadi, suatu fungsi bernilai kompleks $f$ bervariasi terbatas
pada $D$ jika dan hanya jika terdapat fungsi-fungsi terbatas dan
tak-menurun $f_1,\ldots,f_4:\Bbb R\to\Bbb R$ sedemikian sehingga
$f=f_1-f_2+if_3-if_4$ pada $D\cap\dom f$\cmmnt{ (224D)}.

\spheader 224Kd Misalkan $f$ suatu fungsi bernilai kompleks dan $D$
sembarang subhimpunan dari $\Bbb R$. Jika $f$ bervariasi terbatas pada
$D$, maka

\Centerline{$\lim_{x\downarrow a}\Var_{D\cap\ocint{a,x}}(f)
=\lim_{x\uparrow a}\Var_{D\cap\coint{x,a}}(f)=0$}

\noindent untuk setiap $a\in\Bbb R$, dan

\Centerline{$\lim_{a\to-\infty}\Var_{D\cap\ocint{-\infty,a}}(f)
=\lim_{a\to\infty}\Var_{D\cap\coint{a,\infty}}(f)=0$.}

\prooflet{\noindent (Terapkan 224E pada bagian real dan bagian imajiner
dari $f$.)}

\spheader 224Ke Misalkan $f$ suatu fungsi bernilai kompleks yang
bervariasi terbatas pada $[a,b]$, dengan $a<b$. Jika $\dom f$ beririsan
dengan setiap interval $\ocint{a,a+\delta}$ untuk $\delta>0$, maka
 $\lim_{t\in\dom f,t\downarrow a}f(t)$
ada dalam $\Bbb C$. Jika $\dom f$ beririsan dengan
$\coint{b-\delta,b}$ untuk setiap $\delta>0$, maka
$\lim_{t\in\dom f,t\uparrow b}f(t)$
ada dalam $\Bbb C$. \prooflet{(Terapkan 224F pada bagian real
dan bagian imajiner dari $f$.)}

\spheader 224Kf Misalkan $f$, $g$ fungsi-fungsi bernilai kompleks dan $D$ suatu
subhimpunan dari $\Bbb R$. Jika $f$ dan $g$ bervariasi terbatas pada $D$,
demikian pula $f\times g$. \prooflet{(Sebab $f\times g$ dapat dinyatakan
sebagai kombinasi linear dari keempat hasil kali
$\Real f\times\Real g,\ldots,\Imag f\times\Imag g$,
dan 224G dapat diterapkan pada masing-masing hasil kali tersebut.)}

\spheader 224Kg Misalkan $I\subseteq\Bbb R$ suatu interval, dan
$f:I\to\Bbb C$ suatu fungsi bervariasi terbatas. Maka $f$ dapat
didiferensialkan hampir di mana-mana dalam $I$, dan
$\int_I|f'|\le\Var_I(f)$.
\prooflet{(Seperti 224I.)}

\spheader 224Kh Misalkan $f$ dan $g$ fungsi-fungsi bernilai kompleks yang
didefinisikan pada subhimpunan-subhimpunan $\Bbb R$, dan andaikan bahwa
$g$ terintegralkan pada interval $[a,b]$, dengan $a<b$, sedangkan $f$
bervariasi terbatas pada $\ooint{a,b}$ dan didefinisikan hampir di
mana-mana dalam $\ooint{a,b}$. Maka $f\times g$ terintegralkan pada
$[a,b]$, dan

$$\bigl|\int_a^bf\times g\bigr|
\le\bigl(\lim_{x\in\dom f,x\uparrow b}|f(x)|+\Var_{\ooint{a,b}}(f)\bigr)
  \sup_{c\in[a,b]}\bigl|\int_a^cg\bigr|.$$

\prooflet{\noindent (Argumen 224J berlaku nyaris tanpa perubahan.)}

\exercises{
\leader{224X}{Latihan dasar $\pmb{>}$(a)}
%\spheader 224Xa
Tetapkan $f(x)=x^2\sin\Bover1{x^2}$ untuk $x\ne 0$, $f(0)=0$. Tunjukkan
bahwa $f:\Bbb R\to\Bbb R$ terdiferensialkan di setiap titik dan
kontinu seragam, tetapi tidak bervariasi terbatas pada interval tak-sepele
mana pun yang memuat $0$.
%224A

\spheader 224Xb Berikan contoh suatu fungsi nonnegatif
$g:[0,1]\to[0,1]$ yang bervariasi terbatas, tetapi $\sqrt{g}$ tidak
bervariasi terbatas.
%224A

\spheader 224Xc Tunjukkan bahwa jika $f$ sembarang fungsi bernilai real
yang didefinisikan pada suatu subhimpunan dari $\Bbb R$, terdapat fungsi
$\tilde f:\Bbb R\to\Bbb R$ yang memperluas $f$ sedemikian sehingga
$\Var\tilde f=\Var f$. Dalam keadaan apa $\tilde f$ tunggal?
%224D

\spheader 224Xd Misalkan $f:D\to\Bbb R$ suatu fungsi bervariasi terbatas,
dengan $D\subseteq\Bbb R$ suatu himpunan tak kosong. Tunjukkan bahwa jika
$\inf_{x\in
D}|f(x)|>0$, maka $1/f$ bervariasi terbatas.
%224C

\spheader 224Xe Misalkan $f:[a,b]\to\Bbb R$ suatu fungsi kontinu, dengan
$a\le b$ dalam $\Bbb R$. Tunjukkan bahwa jika $c<\Var f$, terdapat
$\delta>0$
sedemikian sehingga $\sum_{i=1}^n|f(a_i)-f(a_{i-1})|\ge c$ kapan pun
$a=a_0\le a_1\le\ldots\le a_n=b$ dan
$\max_{1\le i\le n}a_i-a_{i-1}\le\delta$.
%224C

\spheader 224Xf Misalkan $\sequencen{f_n}$ suatu barisan fungsi bernilai
real, dan tetapkan $f(x)=\lim_{n\to\infty}f_n(x)$ kapan pun limit tersebut
terdefinisi. Tunjukkan bahwa $\Var f\le\liminf_{n\to\infty}\Var f_n$.
%224A

\spheader 224Xg Misalkan $f$ suatu fungsi bernilai real yang
terintegralkan pada interval $[a,b]\subseteq\Bbb R$. Tetapkan
$F(x)=\int_a^xf$ untuk $x\in[a,b]$.
Tunjukkan bahwa $\Var F=\int_a^b|f|$. ({\it Petunjuk\/}: mulailah dengan
memeriksa bahwa $\Var F\le\int|f|$; untuk ketaksamaan sebaliknya, tinjau
terlebih dahulu kasus $f\ge 0$.)
%224D

\spheader 224Xh Tunjukkan bahwa jika $f$ suatu fungsi bernilai real yang
didefinisikan pada himpunan tak kosong $D\subseteq\Bbb R$, maka

\Centerline{$\Var f=\sup\{|\sum_{i=1}^n(-1)^i(f(a_i)-f(a_{i-1}))|:
  a_0\le a_1\le\ldots\le a_n$ dalam $D\}$.}
%224A

\spheader 224Xi Misalkan $f$ suatu fungsi bernilai real yang
terintegralkan pada interval terbatas $[a,b]\subseteq\Bbb R$. Tunjukkan
bahwa

\Centerline{$\int_a^b|f|
=$$\sup\{|\sum_{i=1}^n(-1)^i\int_{a_{i-1}}^{a_i}f|:
  a=a_0\le a_1\le a_2\le\ldots\le a_n=b\}$.}

\noindent({\it Petunjuk\/}: gabungkan 224Xg dan 224Xh.)
%224Xg, 224Xh, 224D

\spheader 224Xj Misalkan $f$ dan $g$ fungsi-fungsi bernilai real yang
didefinisikan pada subhimpunan-subhimpunan $\Bbb R$, dan andaikan bahwa
$g$ terintegralkan pada interval $[a,b]$, dengan $a<b$, sedangkan $f$
bervariasi terbatas pada $\ooint{a,b}$ dan didefinisikan hampir di
mana-mana dalam $\ooint{a,b}$. Tunjukkan bahwa

\Centerline{$|\int_a^bf\times g|
\le(\lim_{x\in\dom f,x\downarrow a}|f(x)|
  +\Var_{\ooint{a,b}}(f))\sup_{c\in[a,b]}|\int_c^bg|$.}
%224J

\spheader 224Xk\dvAnew{2012} Andaikan bahwa $D\subseteq\Bbb R$ dan
$f:D\to\Bbb R$ suatu fungsi. Tunjukkan bahwa $f$ dapat dinyatakan sebagai
selisih dua fungsi tak-menurun jika dan hanya jika
$\Var_{D\cap[a,b]}(f)$ berhingga kapan pun $a\le b$ dalam $D$.
%224D out of order query

\spheader 224Xl\dvAnew{2012}
Andaikan bahwa $D\subseteq\Bbb R$ dan bahwa $f:D\to\Bbb R$ suatu fungsi
kontinu yang bervariasi terbatas. Tunjukkan bahwa $f$ dapat dinyatakan
sebagai selisih dua fungsi kontinu dan tak-menurun.
%224E out of order query

\spheader 224Xm\dvAnew{2013}
Andaikan bahwa $D\subseteq\Bbb R$ dan bahwa $f:D\to\Bbb R$ suatu fungsi
bervariasi terbatas yang kontinu dari kanan, yakni kapan pun $x\in D$ dan
$\epsilon>0$, terdapat $\delta>0$ sedemikian sehingga
$|f(t)-f(x)|\le\epsilon$ untuk setiap $t\in D\cap[x,x+\delta]$.
Tunjukkan bahwa $f$ dapat dinyatakan sebagai selisih dua fungsi
tak-menurun yang kontinu dari kanan.
%224E out of order query

\leader{224Y}{Latihan lanjutan (a)}
%\spheader 224Ya
Tunjukkan bahwa jika $f$ sembarang fungsi bernilai kompleks yang
didefinisikan pada suatu subhimpunan dari $\Bbb R$, terdapat fungsi
$\tilde f:\Bbb R\to\Bbb C$ yang memperluas $f$ sedemikian sehingga
$\Var\tilde f=\Var f$. Dalam keadaan apa $\tilde f$ tunggal?

\spheader 224Yb Misalkan $D$ sembarang subhimpunan tak kosong dari
$\Bbb R$, dan misalkan $\eusm V$ suatu ruang fungsi-fungsi
$f:D\to\Bbb R$ yang bervariasi terbatas. Untuk $f\in\eusm V$, tetapkan

\Centerline{$\|f\|=\sup\{|f(t_0)|+\sum_{i=1}^n|f(t_i)-f(t_{i-1})|:
t_0\le\ldots\le t_n\in D\}$.}

\noindent Tunjukkan bahwa (i) $\|\,\|$ merupakan norma pada $\eusm V$
(ii) $\eusm V$ lengkap terhadap $\|\,\|$ (iii)
$\|f\times g\|\le\|f\|\|g\|$ untuk semua $f$, $g\in\eusm V$, sehingga
$\eusm V$ merupakan suatu aljabar Banach.

\spheader 224Yc Misalkan $f:\Bbb R\to\Bbb R$ suatu fungsi bervariasi
terbatas.
Tunjukkan bahwa terdapat suatu barisan $\sequencen{f_n}$ dari fungsi-fungsi
yang terdiferensialkan
sedemikian sehingga $\lim_{n\to\infty}f_n(x)=f(x)$ untuk setiap
$x\in\Bbb R$, $\lim_{n\to\infty}\int|f_n-f|=0$, dan
$\Var(f_n)\le\Var(f)$ untuk setiap $n\in\Bbb N$. \Hint{mulailah dengan
$f$ yang tak-menurun.}

\spheader 224Yd Untuk setiap himpunan terurut parsial $X$ dan setiap
fungsi $f:X\to\Bbb R$, katakan bahwa $\Var_X(f)=0$ jika $X=\emptyset$,
dan jika tidak,

\Centerline{$\Var_X(f)
=\sup\{\sum_{i=1}^n|f(a_i)-f(a_{i-1})|:a_0,a_1,\ldots,a_n\in X,\,a_0\le
a_1\le\ldots\le a_n\}$.}

\noindent Nyatakan dan buktikan hasil-hasil dalam kerangka ini yang
memperumum 224D dan 224Yb. ({\it Petunjuk-petunjuk\/}: $f$ disebut
`tak-menurun' jika $f(x)\le f(y)$ kapan pun $x\le y$; tafsirkan
$\ocint{-\infty,x}$ sebagai $\{y:y\le x\}$.)

\spheader 224Ye Misalkan $(X,\rho)$ suatu ruang metrik dan
$f:[a,b]\to X$ suatu fungsi, dengan $a\le b$ dalam $\Bbb R$. Tetapkan
\ifnum\stylenumber=12

\Centerline{$\Var_{[a,b]}(f)
=\sup\{\sum_{i=1}^n\rho(f(a_i),f(a_{i-1})):
a\le a_0\le\ldots\le a_n\le b\}$.}

\noindent\else
$\Var_{[a,b]}(f)
=\sup\{\sum_{i=1}^n\rho(f(a_i),f(a_{i-1})):
a\le a_0\le\ldots\le a_n\le b\}$. \fi
(i) Tunjukkan bahwa
$\Var_{[a,b]}(f)=\Var_{[a,c]}(f)+\Var_{[c,b]}(f)$ untuk setiap
$c\in[a,b]$.
(ii) Tunjukkan bahwa jika $\Var_{[a,b]}(f)$ berhingga, maka $f$ kontinu
pada $[a,b]$, kecuali pada terhitung banyak titik.
(iii) Tunjukkan bahwa jika $X$ lengkap dan $\Var_{[a,b]}(f)<\infty$, maka
$\lim_{t\uparrow x}f(t)$ ada untuk setiap $x\in\ocint{a,b}$.
(iv) Tunjukkan bahwa jika $X$ lengkap, maka $\Var_{[a,b]}(f)$ berhingga
jika dan hanya jika $f$ dapat dinyatakan sebagai komposisi $gh$, dengan
$h:[a,b]\to\Bbb R$ tak-menurun dan $g:\Bbb R\to X$ bersifat
$1$-Lipschitz, yakni $\rho(g(c),g(d))\le|c-d|$ untuk semua $c$,
$d\in\Bbb R$.

\spheader 224Yf Misalkan $U$ suatu ruang bernorma dan $a\le b$ dalam
$\Bbb R$. Untuk fungsi-fungsi $f:[a,b]\to U$, definisikan
$\Var_{[a,b]}(f)$ seperti dalam 224Ye, dengan menggunakan metrik baku
$\rho(x,y)=\|x-y\|$ untuk $x$, $y\in U$. (i) Tunjukkan bahwa
$\Var_{[a,b]}(f+g)\le\Var_{[a,b]}(f)+\Var_{[a,b]}(g)$,
$\Var_{[a,b]}(cf)=|c|\Var_{[a,b]}(f)$ untuk semua $f$, $g:[a,b]\to U$
dan semua $c\in\Bbb R$. (ii) Tunjukkan bahwa jika $V$ ruang bernorma lain
dan $T:U\to V$ suatu operator linear terbatas, maka
$\Var_{[a,b]}(Tf)\le\|T\|\Var_{[a,b]}(f)$ untuk setiap $f:[a,b]\to U$.

\spheader 224Yg Misalkan $f:[0,1]\to\Bbb R$ suatu fungsi kontinu. Untuk
$y\in\Bbb R$, tetapkan $h(y)=\#(f^{-1}[\{y\}])$ jika nilai ini berhingga,
dan $\infty$ jika tidak. Tunjukkan bahwa (jika kita mengizinkan $\infty$
sebagai nilai integral) $\Var_{[0,1]}(f)=\int h$. ({\it Petunjuk\/}:
untuk $n\in\Bbb N$, $i<2^n$, tetapkan
$c_{ni}=\sup\{f(x)-f(y):x,y\in[2^{-n}i,2^{-n}(i+1)]\}$,
$h_{ni}(y)=1$ jika $y\in f[\,\coint{2^{-n}i,2^{-n}(i+1)}\,]$, dan $0$
jika tidak. Tunjukkan bahwa $c_{ni}=\int h_{ni}$,
$\lim_{n\to\infty}\sum_{i=0}^{2^n-1}c_{ni}=\Var f$,
$\lim_{n\to\infty}\sum_{i=0}^{2^n-1}h_{ni}=h$.)
(Lihat juga 226Yc.)

\spheader 224Yh Misalkan $\nu$ sembarang ukuran Lebesgue-Stieltjes pada
$\Bbb R$, $I\subseteq \Bbb R$ suatu interval (yang dapat terbuka atau
tertutup, terbatas atau tak terbatas), dan $D\subseteq I$ suatu himpunan
tak kosong. Misalkan $\eusm V$ ruang fungsi-fungsi bervariasi terbatas
dari $D$ ke $\Bbb R$, dan $\|\,\|$ norma 224Yb pada $\eusm V$. Misalkan
$g:D\to\Bbb R$ suatu fungsi sedemikian sehingga
$\int_{[a,b]\cap D}g\,d\nu$ ada kapan pun $a\le b$ dalam $I$, dan
$K=\sup_{a,b\in I,a\le  b}|\int_{[a,b]\cap D}g\,d\nu|$.
Tunjukkan bahwa $|\int_Df\times g\,d\nu|\le K\|f\|$ untuk setiap
$f\in\eusm V$.
%224K

\spheader 224Yi Jelaskan cara menerapkan 224Yh dengan $D=\Bbb N$ untuk
memperoleh Teorema Abel bahwa jika suatu barisan monoton konvergen ke $0$
dan suatu deret mempunyai barisan jumlah parsial yang terbatas, maka deret
hasil kali suku demi suku keduanya dapat dijumlahkan.
%224Yh 224K

\spheader 224Yj Andaikan bahwa $I\subseteq\Bbb R$ suatu interval, dan
bahwa $\sequencen{A_n}$ suatu barisan himpunan yang menutupi $I$.
Misalkan $f:I\to\Bbb R$ kontinu. Tunjukkan bahwa
$\Var f\le\sum_{n=0}^{\infty}\Var_{A_n}f$.
\Hint{cukup buktikan untuk kasus himpunan-himpunan tertutup $A_n$;
gunakan Teorema Baire (4A2Ma).}
}%end of exercises

\endnotes{
\Notesheader{224} Saya telah mengembangkan gagasan-gagasan di atas agak
lebih jauh daripada yang segera kita perlukan; untuk bab ini, cukup meninjau
kasus ketika $D=\dom f=[a,b]$ untuk suatu interval
$[a,b]\subseteq\Bbb R$. Perluasan ke fungsi-fungsi dengan domain tak
beraturan akan berguna dalam Bab 28, dan perluasan ke himpunan tak
beraturan $D$, meskipun tidak penting bagi kita di sini, cukup menarik
-- misalnya, dengan mengambil $D=\Bbb N$, kita memperoleh konsep `barisan
bervariasi terbatas', yang tentu relevan bagi persoalan konvergensi dan
keterjumlahan.

Hasil utama bagian ini tentu saja adalah fakta bahwa suatu fungsi
bervariasi terbatas dapat dinyatakan sebagai selisih fungsi-fungsi
monoton (224D); bahkan, salah satu tujuan konsep ini adalah mencirikan
rentang linear fungsi-fungsi monoton. Hampir semua hasil lain di sini
dapat diperoleh sebagai konsekuensi mudah dari fakta tersebut, seperti
dalam 224E-224G. Dalam 224I dan 224Xg kita melangkah sedikit lebih dalam,
dan bahkan mulai muncul unsur-unsur teori ukuran; di sinilah gagasan-gagasan ini
mulai terhubung dengan pokok sesungguhnya dari bab ini, yang akan
dilanjutkan dalam bagian berikutnya. Hasil lain yang pada dasarnya cukup
mudah, tetapi mengandung benih gagasan-gagasan penting, adalah 224Yg.

Dalam 224Yb saya menyebutkan suatu pengembangan alami dalam analisis
fungsional, dan dalam 224Yd-224Yf %224Yd 224Ye 224Yf
saya mengusulkan generalisasi lebih lanjut yang berjangkauan luas.
}%end of notes

\discrpage
